% ---
% id: receiver/def:compile
% kind: def
% title: "Compilation of a declared element graph"
% summary: >-
%   The compiled instrument matrix is the product of the embedded elements
%   along a linear extension, later elements multiplying from the left,
%   followed by the terminal detector-efficiency layer. This fixes the matrix
%   convention of the whole receiver contract.
% prominence: standard
% status: implemented
% uses: [def:graph, def:instrument]
% mentions: [receiver/detector-imperfections/thm:efficiency]
% assumes: []
% refs: []
% provenance: newly-derived
% implements: certiqs_key.receiver.graph:compile_element_graph
% lean:
% ---

The product below follows the matrix convention the interface records
(\repodoc{spec/modules/receiver/interface.yaml}): a factor acts on column
vectors indexed by the modes, so the element that light meets later stands
further to the left.

\begin{definition}\label{def:compile}
Let $\Grph$ be a well-formed graph; it is in particular acyclic
(condition (E2) of Definition~\ref{def:graph}), so that a linear extension exists. For a
linear extension $\sigma=(e_1,\dots,e_L)$ of $\Grph$ put
\begin{equation}\label{eq:compile}
  \Gmat_{\Grph}^{\sigma}\coloneqq
  \hat\Gmat_{e_L}\,\hat\Gmat_{e_{L-1}}\cdots\hat\Gmat_{e_1}
  \ \in\ \mathbb{C}^{\Modes\times\Modes},
\end{equation}
later elements multiplying from the left. The
\term[The matrix D_eta times the ordered product of a well-formed graph's embedded elements.]{compiled instrument matrix} of the receiver is
$\Gmat\coloneqq D_\eta\,\Gmat^{\sigma}_{\Grph}$, where
$D_\eta=\sum_k\sqrt{\eta_k}\,\Emd{\{k\}}+\Emd{\Modes\setminus\bigcup_k
\Port{k}}$ is the \term[The diagonal matrix D_eta applying each port's detector efficiency after every element.]{terminal detector-efficiency layer}, applied after every
element, and $\eta_k\in[0,1]$ is the declared efficiency of port $k$ (the
effect of that layer on the port Grams is Theorem~\ref{thm:efficiency} of the
detector-imperfections submodule). The \tref{detector ports} $\Port{k}$ and the
port-embedding abbreviation $\Emd{\{k\}}$ written here are those of a
\tref{receiver instrument} tuple: clause (iii) of
Definition~\ref{def:instrument} and its closing notation.
\end{definition}
